\setcounter{section}{13}
\section{Mnożniki Fouriera w \texorpdfstring{$L^2$}{L2}}
\label{sec:dt-multipliers-l2}

W tym rozdziale zachowujemy unitarną normalizację
transformacji Fouriera z definicji \ref{def:fourier-schwartz17}.
Materiał przenoszony z RRC-08 wymaga w tej konwencji
usunięcia czynników $(2\pi)^{-d/2}$ i $(2\pi)^{-d}$ w
tożsamościach norm oraz iloczynów skalarnych.
Por. \atlasref{PDE-08::thm:pde8-multiplier}{RRC-08: mnożniki Fouriera}.


\begin{twierdzenie}[Mnożnik w $L^2$]
\label{thm:dt-multiplier-l2}
Niech $m:\R^d\to\C$ będzie funkcją mierzalną i
istotnie ograniczoną. Dla $f\in L^2(\R^d)$ istnieje
dokładnie jedna klasa $T_mf\in L^2(\R^d)$ spełniająca
$\widehat{T_mf}=m\widehat f$ w sensie $L^2$.
Operator $T_m$ jest liniowy i ograniczony oraz
\[
 \|T_m\|_{L^2\to L^2}=\|m\|_{L^\infty}.
\]
Odwrotnie, jeśli $m$ jest mierzalna i skończona
prawie wszędzie, a ograniczony operator liniowy
$A:L^2\to L^2$ spełnia $\widehat{Af}=m\widehat f$
dla każdego $f\in L^2$ takiego, że
$m\widehat f\in L^2$, to $m\in L^\infty$
oraz $A=T_m$.
\end{twierdzenie}
\begin{proof}
Niech $M=\|m\|_{L^\infty}$. Z definicji istotnego
supremum dla każdego $k\ge1$ zbiór
$\{|m|>M+1/k\}$ ma miarę zero.
Ich przeliczalna suma jest zbiorem $\{|m|>M\}$,
więc $|m|\le M$ prawie wszędzie.
Dla $h\in L^2$ funkcja $mh$ jest więc kwadratowalna
i spełnia $\|mh\|_2\le M\|h\|_2$.
Mnożenie nie zależy od wyboru reprezentantów
klas funkcji $m$ i $h$.
Z twierdzenia~\ref{thm:plancherel22}
transformata Fouriera na $L^2$ jest bijekcją.
Możemy zatem określić
\[
 T_mf=\mathcal F^{-1}(m\widehat f).
\]
Jest to jedyna klasa spełniająca warunek
z tezy, a liniowość wynika z liniowości
wszystkich trzech operatorów w tym złożeniu.
Stosując dwukrotnie twierdzenie Plancherela,
otrzymujemy
\[
 \|T_mf\|_2=\|m\widehat f\|_2
 \le M\|\widehat f\|_2
 =M\|f\|_2.
\]
Stąd $\|T_m\|\le M$.

Gdy $M=0$, funkcja $m$ jest zerowa prawie
wszędzie i $T_m=0$, co daje równość norm.
Załóżmy $M>0$ i ustalmy $0<\varepsilon<M$.
Z definicji istotnego supremum zbiór
$E_\varepsilon=\{\xi:|m(\xi)|>M-\varepsilon\}$
ma dodatnią miarę: gdyby miał miarę zero,
$M-\varepsilon$ byłoby istotnym ograniczeniem
z góry dla $|m|$, wbrew definicji $M$.
Ponieważ $\R^d=\bigcup_{n\ge1}B_n(0)$,
dla pewnego $n$ zbiór
$E=E_\varepsilon\cap B_n(0)$ ma dodatnią
i skończoną miarę. Niech
\[
 h=\frac{\chi_E}{\lambda_d(E)^{1/2}},
 \qquad f=\mathcal F^{-1}h.
\]
Wówczas $\|h\|_2=1$ i
$\|f\|_2=1$. Ponadto
\[
 \|mh\|_2^2
 =\frac1{\lambda_d(E)}\int_E|m(\xi)|^2
 \,\mathrm d\lambda_d(\xi)
 \ge(M-\varepsilon)^2.
\]
W konsekwencji
$\|T_mf\|_2/\|f\|_2=\|mh\|_2\ge M-\varepsilon$.
Przechodząc z $\varepsilon$ do zera,
otrzymujemy $\|T_m\|\ge M$.

Dla implikacji odwrotnej połóżmy $C=\|A\|$.
Jeśli zbiór $\{|m|>C+\varepsilon\}$ dla pewnego
$\varepsilon>0$ ma dodatnią miarę, to z
prawie wszędzie skończoności $m$ wynika, że
z dokładnością do zbioru miary zero jest on
przeliczalną sumą poniższych zbiorów po $n\ge1$.
Dodatnia miara sumy oznacza, że dla pewnego $n$ zbiór
\[
 E=\{\xi:C+\varepsilon<|m(\xi)|\le n\}
       \cap B_n(0)
\]
ma dodatnią miarę. Ma też miarę skończoną.
Wybierając $f=\mathcal F^{-1}\chi_E$,
mamy $m\widehat f\in L^2$ dzięki ograniczeniu
$|m|\le n$ na $E$. Z założenia i Plancherela
wynika wtedy
\[
 \frac{\|Af\|_2}{\|f\|_2}
 =\frac{\|m\chi_E\|_2}{\|\chi_E\|_2}
 \ge C+\varepsilon>C,
\]
co przeczy definicji normy operatorowej.
Zatem dla każdego $\varepsilon>0$ zbiór
$\{|m|>C+\varepsilon\}$ ma miarę zero.
Biorąc przeliczalną sumę dla
$\varepsilon=1/k$, $k\ge1$, dostajemy
$|m|\le C$ prawie wszędzie. Tym samym
$m\in L^\infty$, a warunek na $A$ obowiązuje
już dla wszystkich $f\in L^2$.
Jednoznaczność daje $A=T_m$.
\end{proof}

\begin{wniosek}[Rachunek ograniczonych mnożników]
Dla $m,n\in L^\infty(\R^d)$ i $a,b\in\C$ zachodzi
\[
 T_{am+bn}=aT_m+bT_n,\qquad
 T_mT_n=T_{mn},\qquad T_1=I,\qquad
 T_m^*=T_{\overline m}.
\]
Gwiazdka oznacza operator sprzężony w zespolonej
przestrzeni Hilberta $L^2$.
\end{wniosek}
\begin{proof}
Dla $f\in L^2$ z definicji mnożników mamy
\[
 \begin{aligned}
 \mathcal F(aT_mf+bT_nf)&=(am+bn)\widehat f,\\
 \mathcal F(T_mT_nf)&=m\mathcal F(T_nf)=mn\widehat f,\\
 \mathcal F(T_1f)&=\widehat f.
 \end{aligned}
\]
Wszystkie symbole po prawej stronie są istotnie
ograniczone. Równości operatorów wynikają
z iniektywności transformaty na $L^2$.
Dla $f,g\in L^2$ tożsamość
 z twierdzenia~\ref{thm:plancherel22} daje
\[
 \begin{aligned}
 \langle T_m f,g\rangle_2
 &=\int_{\R^d}m(\xi)\widehat f(\xi)
       \overline{\widehat g(\xi)}\,\mathrm d\lambda_d(\xi)\\
 &=\int_{\R^d}\widehat f(\xi)
       \overline{\overline{m(\xi)}\widehat g(\xi)}
       \,\mathrm d\lambda_d(\xi)\\
 &=\langle f,T_{\overline m}g\rangle_2.
 \end{aligned}
\]
Całki są bezwzględnie zbieżne przez nierówność
Cauchy'ego--Schwarza i istotną ograniczoność $m$.
To dowodzi równości operatorów sprzężonych.
\end{proof}

\begin{przyklad}[Transformata Riesza w $L^2$]
Dla $1\le j\le d$ połóżmy
$m_j(\xi)=-i\xi_j/|\xi|$ dla $\xi\ne0$ i $m_j(0)=0$.
Twierdzenie~\ref{thm:dt-multiplier-l2} definiuje
ograniczony operator $R_j=T_{m_j}:L^2\to L^2$,
zwany transformatą Riesza.
Mamy $|m_j(\xi)|\le1$. Dla każdego
$0<\varepsilon<1$ zbiór
$\{\xi\ne0:|m_j(\xi)|>1-\varepsilon\}$
zawiera otwartą kulę o dodatnim promieniu
wokół $e_j$: funkcja $|\xi_j|/|\xi|$ jest
ciągła w otoczeniu $e_j$ i przyjmuje tam
wartość $1$ w punkcie $e_j$.
Zatem $\|m_j\|_\infty=1$ i w konsekwencji
$\|R_j\|_{L^2\to L^2}=1$, także dla $d=1$.

Ponieważ $\overline{m_j}=-m_j$ oraz dla
$\xi\ne0$ zachodzą
\[
 \sum_{j=1}^d |m_j(\xi)|^2=1,\qquad
 \sum_{j=1}^d m_j(\xi)^2=-1,
\]
z rachunku mnożników otrzymujemy
\[
 R_j^*=-R_j,\qquad
 \sum_{j=1}^d R_j^2=-I,\qquad
 \sum_{j=1}^d\|R_j f\|_2^2=\|f\|_2^2.
\]
Ostatnia tożsamość wynika również bezpośrednio
z Plancherela:
\[
 \sum_{j=1}^d\|R_j f\|_2^2
 =\int_{\R^d}
 \left(\sum_{j=1}^d|m_j(\xi)|^2\right)
 |\widehat f(\xi)|^2\,\mathrm d\lambda_d(\xi)
 =\|f\|_2^2.
\]
Wartość mnożnika w pojedynczym punkcie $\xi=0$
nie zmienia operatora na $L^2$, ponieważ
$\lambda_d(\{0\})=0$. Nie oznacza to jednak,
że można przez ten sam symbol mnożyć dowolną
dystrybucję temperowaną: wzdłuż $\xi=te_j$
symbol przyjmuje wartość $-i$ dla $t>0$
oraz $i$ dla $t<0$. Nie ma więc ciągłego
przedłużenia do zera i nie spełnia warunku
mnożenia dystrybucji z poprzedniej sekcji.
Na przykład dla testu $\varphi$ z $\varphi(0)\ne0$
iloczyn $m_j\varphi$ nie jest ciągły w zerze,
a zatem nie należy do $\mathcal S$.
\end{przyklad}

\begin{uwaga}[Zakres twierdzenia o mnożnikach]
Kryterium istotnej ograniczoności symbolu
dotyczy normy $L^2$ i korzysta z twierdzenia
Plancherela. Dla $p\ne2$ powyższy dowód
nie daje twierdzenia o ograniczoności
na $L^p$. W szczególności nie należy
wyciągać z niego wniosku o ograniczoności
transformat Riesza w normach $L^1$ lub
$L^\infty$. Badanie mnożników na $L^p$
dla $1<p<\infty$ wymaga dalszych narzędzi
analizy harmonicznej; do tej teorii wrócimy po omówieniu operatorów maksymalnych i interpolacji.
\end{uwaga}
